\section{The failure}
\label{sec:failure}

Table~\ref{tab:gate} reproduces the gate's output on the published run, next to
two later values for each count: the first correction, and the value after the
check against object histories described in Section~\ref{sec:second}. The gate
exited zero. Every count in this paper is a number of objects present in the
catalogues concerned, and Appendix~\ref{app:ledger} tracks each disposition
object through all three versions.

\begin{table}[t]
\centering
\small
\caption{Gate output on the published run. Every row agreed between the two
paths, and three rows were wrong. The first correction is the value published
as a fix; the full-history value is the check in Section~\ref{sec:second}.}
\label{tab:gate}
\begin{tabular}{lrrcrr}
\toprule
Defect class & Python & SPARQL & Agree & First correction & Full history \\
\midrule
PhantomEntry              & 22    & 22    & yes & 22 & 22 \\
PhantomEntryOnOrbit       & 1     & 1     & yes & 1 & 1 \\
UndisclosedTrackingLoss   & 1{,}104 & 1{,}104 & yes & \textbf{1{,}094} & 1{,}094 \\
DispositionDisagreement   & 932   & 932   & yes & \textbf{261} & \textbf{220} \\
CoverageGap               & 900   & 900   & yes & \textbf{622} & 622 \\
UnnumberedObject          & 605   & 605   & yes & 605 & 605 \\
IdentifierCollision       & 3     & 3     & yes & 3 & 3 \\
\bottomrule
\end{tabular}
\end{table}

The three errors have different characters, and only the first is the
common-mode failure this paper is about.

\paragraph{Disposition disagreement: 932, then 261, then 220.}
The count is the number of objects on which the two catalogues contradict each
other about whether the object is still in orbit. The 932 comprised 163 objects
GCAT records as gone while CelesTrak gives no decay date, and 769 objects
CelesTrak records as decayed while GCAT does not record them as gone. Producing
the count requires reading GCAT's status codes, and the pre-correction rule was
simple: a status in a set named \texttt{GONE} meant the object was gone, and
every other status meant it was still in orbit. That default did most of the
damage. Of the 769, 640 were objects whose GCAT record ends in a code that marks
a transition rather than a fate, such as docking, attachment, transfer or
renaming, and 500 of those were docking or attachment specifically. Another 57
ended in an explosion or a collision. The first correction removed 998
transition-coded objects from the comparison, 640 of which had been counted as
disagreements, and reclassified explosion and collision as destruction, giving
261. The ledger in Appendix~\ref{app:ledger} accounts for every object that
entered or left. Section~\ref{sec:second} shows that 261 was itself wrong.

\paragraph{Coverage gap: 900 against 622.}
The pipeline checked three of the four main catalogue files the source
publishes. The fourth was listed on the same index page and never fetched. It
holds 354 objects, 278 of which were being counted as absent from a catalogue
that contains them.

\paragraph{Undisclosed tracking loss: 1,104 against 1,094.}
Ten of the 1,104 objects do carry the CelesTrak field whose absence the count
is about. Correcting this strengthened the finding. The field is maintained and
set on 1,292 objects across the catalogue, ten of them among the 1,104, so the
claim moved from ``this register cannot record tracking loss'' to ``this
register records tracking loss and does not record it for 1,094 objects.''

None of the three errors is visible to a redundancy check, and only the first is
common-mode in the strict sense. What unites them is that each is an error about
what the source data means or contains, and the gate compares only what two
implementations derive from it.
